% !TEX program = lualatex
\documentclass[11pt]{ltjarticle}
\usepackage{luatexja}
\usepackage[margin=20truemm]{geometry}
\usepackage{titlesec}
\usepackage{pdfpages}
\usepackage{amssymb}
\usepackage{amsmath}
\usepackage{graphicx}
\usepackage{float}
\usepackage{url}
\usepackage{xcolor}
\usepackage[hidelinks]{hyperref}
\usepackage{listings}
\usepackage{tcolorbox}

% 句読点の置換
\catcode`、=\active
\protected\def、{，}
\catcode`。=\active
\protected\def。{．}

% セクション書体の設定
\titleformat*{\section}{\large\bfseries}
\renewcommand{\baselinestretch}{1.2}

% コード表示の設定
\lstset{
  language={bash},
  basicstyle=\ttfamily\small,
  breaklines=true,
  frame=single,
  numberstyle=\tiny,
  showstringspaces=false,
  tabsize=2,
  aboveskip=\medskipamount,
  belowskip=\medskipamount,
  backgroundcolor=\color{lightgray!20}
}

\title{Windows環境でのLaTeX環境構築ガイド}
\author{経験者からのアドバイス}
\date{\today}

\begin{document}

\maketitle

\tableofcontents

\newpage

\section{はじめに}

本ドキュメントは，Windows環境でLaTeXを使用してレポートを作成する際に必要な環境構築手順を説明します．

このガイドは実際の運用経験に基づいており，以下を実現することを目標としています：

\begin{itemize}
  \item 日本語に対応した高品質なPDF出力
  \item 複数のプリアンブルファイルと図表の効率的な管理
  \item 自動化されたコンパイルプロセス
  \item エラー処理と変更検出による効率化
\end{itemize}

\section{推奨構成}

本ガイドで推奨する構成は以下の通りです：

\begin{center}
\begin{tabular}{ll}
\hline
\textbf{項目} & \textbf{詳細} \\
\hline
LaTex配布版 & TeX Live 2024以降 \\
コンパイラ & LuaLaTeX \\
エディタ & Visual Studio Code \\
Python & 3.9以上（オプション：自動化用） \\
OS & Windows 10 / 11 \\
\hline
\end{tabular}
\end{center}

\section{最小限のパッケージ構成}

\subsection{実際の運用に必要なパッケージ}

当テンプレートを使用する場合，以下のパッケージが\textbf{最小限必須}です：

\begin{itemize}
  \item \textbf{luatexja}：日本語組版（LuaTeXベース）
  \item \textbf{geometry}：ページレイアウト設定
  \item \textbf{titlesec}：セクション見出しのカスタマイズ
  \item \textbf{amsmath, amssymb}：数式環境と数学記号
  \item \textbf{graphicx}：画像の挿入
  \item \textbf{float}：図表の位置固定 [H] オプション
  \item \textbf{hyperref}：ハイパーリンク（参照・URL）
  \item \textbf{siunitx}：単位と数値の表示
  \item \textbf{pgfplots, tikz}：グラフ描画
\end{itemize}

\subsection{オプションパッケージ（分野別）}

分野に応じて以下を追加します：

\begin{itemize}
  \item \textbf{化学分野}：mhchem（化学式），chemfig（構造式）
  \item \textbf{電気・電子分野}：circuitikz（回路図）
  \item \textbf{一般的な活用}：pdfpages（PDF挿入），url（URL表示），xcolor（色使用）
\end{itemize}

\subsection{TeX Liveの最小インストール戦略}

TeX Liveのインストール時に以下の選択肢があります：

\begin{enumerate}
  \item \textbf{フルインストール（推奨）}
    \begin{itemize}
      \item すべてのパッケージをインストール
      \item サイズ：約6〜7GB
      \item 利点：追加パッケージ不要，スムーズな運用
      \item 欠点：ディスク容量が大きい
    \end{itemize}

  \item \textbf{カスタムインストール（上級者向け）}
    \begin{itemize}
      \item 必要なパッケージのみを選択
      \item サイズ：約2〜3GB（最小構成）
      \item 利点：ディスク容量削減
      \item 欠点：後から追加パッケージを導入する手間
    \end{itemize}
\end{enumerate}

\textbf{推奨}：初めてのインストールは\textbf{フルインストール}を選択することをお勧めします．これにより，パッケージの依存関係に悩まされず，安定した運用が可能になります．

\subsection{インストール後のパッケージ確認}

TeX Liveをインストール後，以下のコマンドで必須パッケージがインストールされているか確認できます：

\begin{lstlisting}[language=bash]
# Windows PowerShellで実行
& "C:\texlive\2024\bin\windows\tlmgr.bat" list --installed | `
  Select-String "luatexja|geometry|titlesec|amsmath|amssymb"
\end{lstlisting}

各パッケージが表示されれば，インストール完了です．

\section{インストール手順}

\subsection{1. TeX Liveのインストール}

TeX Liveは，LaTeXの完全な配布版です．日本語対応や多数のパッケージが含まれています．

\subsubsection{ダウンロード}

以下の方法でインストーラを取得します：

\begin{enumerate}
  \item \url{https://www.tug.org/texlive/acquire-netinstall.html} にアクセス
  \item \texttt{install-tl-windows.exe} をダウンロード
  \item インストーラを実行
\end{enumerate}

\subsubsection{インストール設定}

実行時に表示されるダイアログで以下を確認します：

\begin{itemize}
  \item \textbf{言語選択}：Japanese を含める
  \item \textbf{パッケージスキーム}：Full（全パッケージ）を推奨
  \item \textbf{インストール先}：\texttt{C:\textbackslash texlive\textbackslash 2024}など
  \item \textbf{オプション}：
    \begin{itemize}
      \item Ghostscript：有効にする（PDF処理用）
      \item Perl：有効にする（スクリプト実行用）
    \end{itemize}
\end{itemize}

インストール完了には時間がかかります（目安：30分～1時間）．

\subsubsection{インストール後の確認}

PowerShellで以下を実行し，インストール確認します：

\begin{lstlisting}[language=bash]
lualatex --version
\end{lstlisting}

バージョン情報が表示されればインストール成功です．

\subsection{2. Visual Studio Codeのセットアップ}

VS Codeは軽量で拡張機能が充実しており，LaTeX編集に最適です．

\subsubsection{インストール}

\begin{enumerate}
  \item \url{https://code.visualstudio.com/} からVS Codeをダウンロード
  \item インストーラを実行
  \item デフォルト設定でインストール完了
\end{enumerate}

\subsubsection{必須拡張機能}

VS Code内で以下の拡張機能をインストールします（拡張機能パネルで検索）：

\begin{enumerate}
  \item \textbf{LaTeX Workshop}（James Yu）
    \begin{itemize}
      \item LaTeX編集の主要拡張機能
      \item シンタックスハイライト，コンパイル，プレビュー
    \end{itemize}
  \item \textbf{LaTeX language support}（Joseph Benden）
    \begin{itemize}
      \item 追加の言語サポート
    \end{itemize}
  \item \textbf{Japanese Language Pack}
    \begin{itemize}
      \item VS Codeのメニューを日本語化
    \end{itemize}
  \item \textbf{Code Spell Checker}
    \begin{itemize}
      \item スペル・文法チェック（オプション）
    \end{itemize}
\end{enumerate}

\subsubsection{LaTeX Workshopの設定}

VS Codeの \texttt{settings.json} に以下を追加します：

\begin{lstlisting}[basicstyle=\ttfamily\tiny]
"latex-workshop.latex.tools": [
  {
    "name": "lualatex",
    "command": "lualatex",
    "args": ["-interaction=nonstopmode", "-halt-on-error", "%DOC%"]
  }
],
"latex-workshop.latex.recipes": [
  {
    "name": "lualatex",
    "tools": ["lualatex"]
  }
],
"latex-workshop.view.pdf.viewer": "external"
\end{lstlisting}

\subsection{3. Pythonのセットアップ（オプション）}

自動コンパイルや変更検出を使用する場合，Pythonが必要です．

\subsubsection{インストール}

\begin{enumerate}
  \item \url{https://www.python.org/} からPython 3.9以上をダウンロード
  \item インストーラを実行
  \item \textbf{重要}：「Add Python to PATH」をチェック
\end{enumerate}

\subsubsection{インストール確認}

\begin{lstlisting}[language=bash]
python --version
pip --version
\end{lstlisting}

\section{ディレクトリ構造}

レポートプロジェクトの推奨ディレクトリ構造は以下の通りです：

\begin{lstlisting}[language=bash]
my_report/
├── main.tex          # メインファイル
├── macros.tex        # マクロ定義（タイトル，著者等）
├── figure.tex        # グラフ設定
├── frontcover.tex    # 表紙
├── purpose.tex       # 目的
├── principle.tex     # 原理
├── method.tex        # 実験方法
├── result.tex        # 結果
├── consideration.tex # 考察
├── assignment.tex    # 課題
├── conclusion.tex    # 結論
├── references.tex    # 参考文献
└── data/             # グラフデータ（CSVなど）
    ├── temperature.csv
    └── frequency.csv
\end{lstlisting}

\subsection{ファイルの役割}

\begin{itemize}
  \item \textbf{main.tex}：全体の統合ファイル，パッケージ読み込み，本文構成
  \item \textbf{macros.tex}：レポートのメタ情報（タイトル，学籍番号，著者）とショートカット
  \item \textbf{figure.tex}：グラフの定義，pgfplotsの設定
  \item \textbf{各セクション*.tex}：本文内容（各セクションの詳細）
\end{itemize}

\section{基本的なLaTeX構文}

\subsection{セクション構成}

\begin{lstlisting}[language=bash, basicstyle=\ttfamily\tiny]
\section{セクション名}      % 第1レベル
\subsection{サブセクション}  % 第2レベル
\subsubsection{副題}        % 第3レベル
\end{lstlisting}

\subsection{箇条書き}

\begin{lstlisting}[language=bash, basicstyle=\ttfamily\tiny]
% 番号なし
\begin{itemize}
  \item 項目1
  \item 項目2
\end{itemize}

% 番号あり
\begin{enumerate}
  \item 項目1
  \item 項目2
\end{enumerate}
\end{lstlisting}

\subsection{数式}

\begin{lstlisting}[language=bash, basicstyle=\ttfamily\tiny]
% インライン数式
$E = mc^2$

% 別行立ち数式
\begin{equation}
  y = ax + b
  \label{eq:linear}
\end{equation}

% 参照
式(\ref{eq:linear})より ...
\end{lstlisting}

\subsection{図表の挿入}

\begin{lstlisting}[language=bash, basicstyle=\ttfamily\tiny]
% 図の挿入
\begin{figure}[H]
  \centering
  \includegraphics[width=8cm]{data/graph.png}
  \caption{グラフのタイトル}
  \label{fig:graph}
\end{figure}

% 参照
\figref{fig:graph}に示す通り ...

% 表の挿入
\begin{table}[H]
  \centering
  \begin{tabular}{|c|c|}
    \hline
    項目1 & 項目2 \\
    \hline
    値1 & 値2 \\
    \hline
  \end{tabular}
  \caption{表のタイトル}
  \label{tab:table}
\end{table}
\end{lstlisting}

\subsection{参考文献}

\begin{lstlisting}[language=bash, basicstyle=\ttfamily\tiny]
% 文中での引用
\cite{参考文献ID}

% references.tex での定義
\begin{thebibliography}{99}
  \bibitem{参考文献ID}
  著者名, 「タイトル」, 発行元, 2024.
\end{thebibliography}
\end{lstlisting}

\section{よくあるパッケージ}

\subsection{化学分野}

\begin{itemize}
  \item \textbf{mhchem}：化学式表示
    \begin{lstlisting}[language=bash, basicstyle=\ttfamily\small]
\usepackage[version=4]{mhchem}
\ce{H2O}, \ce{CO2 + H2O -> H2CO3}
    \end{lstlisting}

  \item \textbf{chemfig}：化学構造式
    \begin{lstlisting}[language=bash, basicstyle=\ttfamily\small]
\usepackage{chemfig}
\chemfig{...}  % 化学構造式を描画
    \end{lstlisting}
\end{itemize}

\subsection{電気・電子分野}

\begin{itemize}
  \item \textbf{circuitikz}：回路図
    \begin{lstlisting}[language=bash, basicstyle=\ttfamily\small]
\usepackage{circuitikz}
\begin{circuitikz}
  \draw (0,0) to[R] (2,0);
\end{circuitikz}
    \end{lstlisting}

  \item \textbf{siunitx}：単位と数値の表示（既にmain.texに含まれている）
    \begin{lstlisting}[language=bash, basicstyle=\ttfamily\small]
\usepackage{siunitx}
\SI{10}{\kilo\ohm}, \SI{1.5}{\micro\farad}
    \end{lstlisting}
\end{itemize}

\subsection{グラフ描画}

\begin{itemize}
  \item \textbf{pgfplots}：科学技術計算グラフ（既にfigure.texに含まれている）
  \item \textbf{tikz}：一般的な図形描画（既にfigure.texに含まれている）
\end{itemize}

\section{コンパイルと実行}

\subsection{VS Codeでの実行}

\begin{enumerate}
  \item VS Codeで \texttt{main.tex} を開く
  \item 左サイドバーの「TeX」アイコンをクリック
  \item 「Build LaTeX project」をクリック
  \item エラーなく完了すると \texttt{main.pdf} が生成される
\end{enumerate}

\subsection{コマンドラインでの実行}

\begin{lstlisting}[language=bash]
cd path/to/report
lualatex -interaction=nonstopmode -halt-on-error main.tex
\end{lstlisting}

\subsection{Pythonスクリプトでの自動化}

\texttt{compile.py} を使用すると，複数のレポートを効率的に管理できます：

\begin{lstlisting}[language=bash]
# config.json に報告書を登録
python compile.py report_name

# すべてのレポートをコンパイル
python compile.py
\end{lstlisting}

\texttt{compile.py} の機能：

\begin{itemize}
  \item 変更ファイルの自動検出
  \item スキップ（変更がない場合）
  \item エラー分析と表示
  \item 補助ファイルの自動クリーンアップ
  \item PDFの自動リネーム（macros.texのメタ情報から）
\end{itemize}

\section{最小限構成での実運用}

\subsection{最小限main.tex テンプレート}

以下は，最小限のパッケージで動作するmain.texです：

\begin{lstlisting}[basicstyle=\ttfamily\small]
% !TEX program = lualatex
\documentclass[11pt]{ltjarticle}
\usepackage{luatexja}
\usepackage[margin=20truemm]{geometry}
\usepackage{titlesec}
\usepackage{amsmath}
\usepackage{amssymb}
\usepackage{graphicx}
\usepackage{float}
\usepackage[hidelinks]{hyperref}
\usepackage{siunitx}

% 句読点の置換
\catcode`、=\active
\protected\def、{，}
\catcode`。=\active
\protected\def。{．}

% セクション書体の設定
\titleformat*{\section}{\large\bfseries}

\title{レポートタイトル}
\author{著者名}
\date{\today}

\begin{document}
\maketitle
\section{はじめに}
\end{document}
\end{lstlisting}

このテンプレートで以下が可能です：

\begin{itemize}
  \item \textbf{日本語サポート}：luatexjaにより完全な日本語出力
  \item \textbf{ページレイアウト}：20mmの余白設定
  \item \textbf{数式環境}：amsmath, amssymbで数学記号を完全サポート
  \item \textbf{画像挿入}：graphicx, floatで図表管理
  \item \textbf{ハイパーリンク}：hyperrefで参照とURLをクリック可能に
  \item \textbf{単位表示}：siunitxで標準化された単位表示
\end{itemize}

\subsection{グラフが不要な場合}

pgfplotsやtikzを使用しない場合，figure.texファイルは最小化できます：

\begin{lstlisting}[basicstyle=\ttfamily\small]
% figure.tex（最小版）
% グラフを使用しない場合は，このファイルを削除するか
% main.tex の % ===== グラフ設定ファイル =====
% グラフに関するすべての設定とマクロをここに記述
% main.tex から \input{figure_config.tex} で読み込まれます

% ---- pgfplots/tikz 関連パッケージ ----
\usepackage{pgfplots}
\usepackage{pgfplotstable}
\usepackage{tikz}
\usepgfplotslibrary{fillbetween}

% ---- グラフマーカーの定義 ----
\pgfplotscreateplotcyclelist{reportcyclelist}{%
    {black, mark=o, mark size=3.2pt, line width=0.9pt, mark options={draw=black, fill=none}},%
    {black, mark=square, mark size=3.2pt, line width=0.9pt, mark options={draw=black, fill=none}},%
    {black, mark=triangle, mark size=3.2pt, line width=0.9pt, mark options={draw=black, fill=none}},%
    {black, mark=diamond, mark size=3.2pt, line width=0.9pt, mark options={draw=black, fill=none}},%
    {black, mark=x, mark size=3.2pt, line width=0.9pt}%
}

% ---- グラフ全体の設定 ----
\pgfplotsset{
    compat=1.18,
    every axis/.style={
        width=10cm,
        height=10cm,
        axis lines=box,
        axis line style={black, line width=1.0pt},
        tick style={black, line width=0.8pt},
        tick align=inside,
        tick label style={font=\small},
        label style={font=\small},
        title style={font=\small\bfseries},
        grid=none,
        minor tick num=1,
        scaled ticks=false,
        legend cell align=left,
        legend style={
            draw=black,
            fill=white,
            fill opacity=0.85,
            text opacity=1,
            font=\small,
            rounded corners=1pt,
        },
        cycle list name=reportcyclelist,
    },
    plot area/.style={width=10cm, height=10cm},
    plot square/.style={width=10cm, height=10cm},
    plot logx/.style={xmode=log},
    plot logy/.style={ymode=log},
    plot loglog/.style={xmode=log, ymode=log},
    plot no grid/.style={grid=none},
    fit line/.style={black, no marks, line width=1.1pt, on layer=axis foreground},
    scatter points/.style={only marks, mark=o, mark size=3.2pt, mark options={draw=black, fill=none}, on layer=axis background},
    line series/.style={line width=1.0pt},
    area series/.style={fill opacity=0.18, on layer=axis background},
}

% ========================================
% 補助マクロ
% ========================================

\newcommand{\PlotTable}[4][]{%
    \addplot[scatter points, #1] table [x=#2, y=#3, col sep=comma] {#4};%
}

\newcommand{\PlotFitLine}[4]{%
    % #1=slope, #2=intercept, #3=xmin, #4=xmax
    \addplot[fit line] coordinates {(#3,{#1*(#3)+#2}) (#4,{#1*(#4)+#2})};%
}

% ========================================
% グラフ定義（tikzpicture で描画）
% 以下、各グラフをマクロ化して定義
% ========================================

% ---- 図1: 基本プロット ----
\newcommand{\GraphBasicPlot}{%
    \begin{tikzpicture}
    \begin{axis}[
        plot area,
        title={電圧測定結果},
        xlabel={時間 [s]},
        ylabel={電圧 [V]},
        legend style={at={(0.98,0.98)}, anchor=north east},
    ]
        \addplot[scatter points] coordinates {
            (0, 0)
            (1, 1.2)
            (2, 2.1)
            (3, 3.0)
            (4, 3.8)
            (5, 4.9)
        };
        \addlegendentry{測定値}

        \addplot[fit line] coordinates {
            (0, 0)
            (5, 5.1)
        };
        \addlegendentry{近似直線}
    \end{axis}
    \end{tikzpicture}
}

% ---- 図2: 温度特性 ----
\newcommand{\GraphTemperatureResponse}{%
    \begin{tikzpicture}
    \begin{axis}[
        plot area,
        title={抵抗の温度特性},
        xlabel={温度 [$^\circ$C]},
        ylabel={抵抗 [$\Omega$]},
        legend style={at={(0.98,0.02)}, anchor=south east},
    ]
        \addplot[scatter points] table [x=temperature, y=resistance, col sep=comma] {data/temperature_response.csv};
        \addlegendentry{実験値}

        \addplot[fit line] coordinates {
            (0, 1000)
            (100, 500)
        };
        \addlegendentry{近似直線}
    \end{axis}
    \end{tikzpicture}
}

% ---- 図3: 周波数応答 ----
\newcommand{\GraphFrequencyResponse}{%
    \begin{tikzpicture}
    \begin{axis}[
        plot area,
        title={周波数応答（ゲイン）},
        xlabel={周波数 [Hz]},
        ylabel={ゲイン},
        xmode=log,
        legend style={at={(0.98,0.02)}, anchor=south east},
    ]
        \addplot[scatter points] table [x=frequency, y=magnitude, col sep=comma] {data/frequency_response.csv};
        \addlegendentry{測定値}

        \addplot[fit line] coordinates {
            (1, 0.95)
            (10000, 0.15)
        };
        \addlegendentry{近似曲線}
    \end{axis}
    \end{tikzpicture}
}
 をコメントアウト
\end{lstlisting}

\subsection{分野別の最小構成例}

\subsubsection{化学分野}

\begin{lstlisting}[basicstyle=\ttfamily\small]
% 化学式が必要な場合のみ追加
\usepackage[version=4]{mhchem}

% 使用例
\ce{H2O}, \ce{CO2 + H2O -> H2CO3}
\end{lstlisting}

\subsubsection{電気・電子分野}

\begin{lstlisting}[basicstyle=\ttfamily\small]
% 回路図が必要な場合のみ追加
\usepackage{circuitikz}

% 使用例
\begin{circuitikz}
  \draw (0,0) to[R] (2,0);
\end{circuitikz}
\end{lstlisting}

\section{トラブルシューティング}

\subsection{文字化けが発生する}

\textbf{症状}：日本語が正しく表示されない

\textbf{対策}：

\begin{enumerate}
  \item \texttt{main.tex} の先頭を確認：
    \begin{lstlisting}[language=bash, basicstyle=\ttfamily\small]
% !TEX program = lualatex
\documentclass[11pt]{ltjarticle}
\usepackage{luatexja}
    \end{lstlisting}
  \item ファイルのエンコーディングがUTF-8であることを確認
  \item VS Codeの右下で「UTF-8」が表示されていることを確認
\end{enumerate}

\subsection{パッケージが見つからない}

\textbf{症状}：\texttt{Package not found} エラー

\textbf{対策}：

\begin{enumerate}
  \item TeX Liveを完全インストール（Full）しているか確認
  \item パッケージ名の綴りを確認
  \item 必要に応じてTeX Liveを更新：
    \begin{lstlisting}[language=bash]
tlmgr update --self --all
    \end{lstlisting}
\end{enumerate}

\subsection{グラフが表示されない}

\textbf{症状}：グラフやCSVデータが反映されない

\textbf{対策}：

\begin{enumerate}
  \item \texttt{data/} ディレクトリが正しい位置にあるか確認
  \item CSVファイルが正しい形式か確認（カンマ区切り）
  \item \texttt{figure.tex} のパスが相対パスになっているか確認
\end{enumerate}

\subsection{コンパイルが異常に遅い}

\textbf{症状}：LuaLaTeXの実行に時間がかかる

\textbf{対策}：

\begin{enumerate}
  \item 不要なパッケージを \texttt{main.tex} から削除
  \item \texttt{pgfplots} の図が多くないか確認
  \item キャッシュをクリア：
    \begin{lstlisting}[language=bash]
rm -f *.aux *.log *.pdf *.out
    \end{lstlisting}
\end{enumerate}

\section{応用：マクロの活用}

\subsection{カスタムコマンド定義}

\texttt{macros.tex} にカスタムコマンドを定義することで，書き込みを効率化できます：

\begin{lstlisting}[language=bash, basicstyle=\ttfamily\small]
% 微分演算子
\newcommand{\diff}{\mathrm{d}}

% 全微分
\newcommand{\D}[1]{\frac{\diff#1}{\diff x}}

% 二重括弧
\newcommand{\norm}[1]{\left\lVert#1\right\rVert}

% 部分文字
\newcommand{\mysubA}{実験グループA}
\end{lstlisting}

使用例：

\begin{lstlisting}[language=bash, basicstyle=\ttfamily\small]
本実験では\mysubA{}の結果を分析した．
導関数は $\D{y}$ で表される．
\end{lstlisting}

\subsection{条件分岐}

異なるバージョンのレポートを管理する場合：

\begin{lstlisting}[language=bash, basicstyle=\ttfamily\small]
% macros.tex で定義
\newif\ifEnglish
\Englishfalse  % false: 日本語, true: 英語

% main.tex で使用
\ifEnglish
  \section{Introduction}
\else
  \section{はじめに}
\fi
\end{lstlisting}

\section{フロー図：環境構築手順}

以下は環境構築全体のフロー図です：

\begin{enumerate}
  \item \textbf{準備フェーズ}
    \begin{enumerate}
      \item TeX Liveをダウンロード
      \item VS Codeをダウンロード
      \item Python（オプション）をダウンロード
    \end{enumerate}

  \item \textbf{インストールフェーズ}
    \begin{enumerate}
      \item TeX Liveをインストール（フルインストール）
      \item VS Codeをインストール
      \item Pythonをインストール（PATH設定有効）
    \end{enumerate}

  \item \textbf{設定フェーズ}
    \begin{enumerate}
      \item VS Code拡張機能をインストール
      \item \texttt{settings.json} を設定
      \item テストファイルでコンパイルテスト
    \end{enumerate}

  \item \textbf{運用フェーズ}
    \begin{enumerate}
      \item テンプレートをコピー
      \item \texttt{macros.tex} を編集
      \item VS Codeで編集・コンパイル
    \end{enumerate}
\end{enumerate}

\section{インストール済みパッケージの確認と削減}

\subsection{インストール済みパッケージの一覧表示}

TeX Liveにインストールされているパッケージ一覧を表示します：

\begin{lstlisting}[language=bash]
# PowerShellで実行
& "C:\texlive\2024\bin\windows\tlmgr.bat" list --only-installed > installed_packages.txt
\end{lstlisting}

\subsection{不要なパッケージの削除}

インストール済みパッケージが不要な場合は削除できます：

\begin{lstlisting}[language=bash]
# 個別パッケージの削除
& "C:\texlive\2024\bin\windows\tlmgr.bat" remove package_name

# 例：汎用計画言語パッケージの削除
& "C:\texlive\2024\bin\windows\tlmgr.bat" remove metapost
\end{lstlisting}

\textbf{注意}：パッケージの削除時は，依存関係を確認してから実行してください．

\subsection{ディスク容量の見積}

\begin{itemize}
  \item \textbf{TeX Live フルインストール}：約6〜7GB
  \item \textbf{最小限の構成}（luatexja + 基本パッケージ）：約2GB
  \item \textbf{当テンプレート運用に必要}：約3GB
\end{itemize}

\subsection{Windows での容量確認}

\begin{lstlisting}[language=bash]
# PowerShell で TeX Live フォルダのサイズ確認
(Get-ChildItem "C:\texlive" -Recurse |
  Measure-Object -Property Length -Sum).Sum / 1GB
\end{lstlisting}

\section{参考リンク}

\begin{itemize}
  \item TeX Live公式：\url{https://www.tug.org/texlive/}
  \item LaTeX Workshop：\url{https://github.com/James-Yu/LaTeX-Workshop}
  \item 日本語LaTeX環境：\url{https://www.latex.or.jp/}
  \item pgfplotsマニュアル：\url{https://pgfplots.sourceforge.io/}
\end{itemize}

\section{まとめ}

本ドキュメントで紹介した環境構築により，以下を実現できます：

\begin{itemize}
  \item \textbf{効率的な執筆}：テンプレートとマクロによる統一化
  \item \textbf{高品質な出力}：LuaLaTeXと最新パッケージによる最適化
  \item \textbf{自動化}：Pythonスクリプトによる管理
  \item \textbf{保守性}：段階的なファイル分割で可読性向上
\end{itemize}

環境構築後は，テンプレートを活用して効率的なレポート作成を進めてください．

\end{document}
